\documentclass[11pt,a4paper]{article}
\usepackage[utf8]{inputenc}
\usepackage[T1]{fontenc}
\usepackage[ngerman]{babel}
\usepackage[margin=2.4cm]{geometry}
\usepackage{booktabs}
\usepackage{longtable}
\usepackage{xcolor}
\usepackage[colorlinks=true,urlcolor=blue!60!black,linkcolor=black]{hyperref}
\usepackage{parskip}
\usepackage{enumitem}
\setlist{nosep,leftmargin=*}

\newcommand{\code}[1]{\texttt{\small #1}}
\newcommand{\ok}{\textcolor{green!45!black}{\textbf{erledigt}}}
\newcommand{\offen}{\textcolor{orange!75!black}{\textbf{offen}}}
\newcommand{\upstream}{\textcolor{blue!60!black}{\textbf{upstream}}}

\title{AWI-ESM3-4-2-veg-HR, CMIP7 piControl\\[2mm]
\large Was sich seit \code{cli112} geaendert hat, und Antwort auf deine letzten drei Runden}
\author{Jan Streffing, AWI}
\date{17. August 2026}

\begin{document}
% babel-ngerman macht " aktiv, was in Code-Beispielen zuschlaegt:
% operation="instant" wurde als operation=ïnstant" gesetzt. Die Shorthands
% werden erst bei \begin{document} scharfgeschaltet, also muss das hierhin.
\shorthandoff{"}
\maketitle

\section{Kurzfassung}

Seit \code{cli112} sind drei Runden von dir eingegangen: der Koordinaten-Check, die
Liste mit den acht Punkten, und dein Plugin. Zusammen mit dem, was das Plugin selbst
noch gefunden hat, waren das rund zwanzig Befunde. Umfang seither: 15 Commits,
26 Dateien, +1197 / -209 Zeilen.

Der Zwischenstand an harten Zahlen, gemessen an \code{cli114}:

\begin{center}
\begin{tabular}{lrr}
\toprule
& \textbf{cli113} & \textbf{cli114} \\
\midrule
HIGH   & 69  & 1 \\
MEDIUM & 288 & 284 \\
\bottomrule
\end{tabular}
\end{center}

Das verbleibende HIGH ist weiterhin \code{basin} und wartet auf
\code{ESGF/cc-plugin-wcrp\#67}. \code{cli115} laeuft gerade und enthaelt alles aus
diesem Dokument; die Gegenprobe schicken wir nach.

\section{Dein Plugin}

\code{cc-plugin-aicc} laeuft ab sofort in jedem unserer Laeufe mit, als dritte Suite
neben \code{cf} und \code{wcrp\_cmip7}. Das war die wirksamste einzelne Aenderung
dieser Runde: er sieht eine ganze Klasse, die die beiden anderen nicht anfassen.
41~Dateien mit \code{latitude(latitude)} statt \code{lat(lat)} sind an \code{cf} und
\code{wcrp\_cmip7} glatt vorbeigelaufen.

Ein Hinweis dazu, den du vermutlich wissen willst. Ohne gesetztes
\code{CMIP7\_TABLES\_PATH} wirft das Plugin eine \code{FileNotFoundError} aus
\code{ComplianceChecker.run\_checker} heraus, und damit stirbt der gesamte
\code{cchecker}-Aufruf:

\begin{center}
\begin{tabular}{lll}
\toprule
& \textbf{Exit} & \textbf{Report} \\
\midrule
ohne \code{CMIP7\_TABLES\_PATH} & 1 & gar keiner \\
mit \code{CMIP7\_TABLES\_PATH}  & 1 & alle drei Suiten \\
\bottomrule
\end{tabular}
\end{center}

Es fehlt also nicht nur das aicc-Ergebnis, sondern auch das von \code{cf} und
\code{wcrp\_cmip7} fuer dieselbe Datei. Bei uns waere das lautlos gewesen: pycmor
protokolliert, dass kein JSON entstand, und laeuft weiter. Eine Datei ohne Sidecar
wird nirgends gezaehlt, der Lauf haette also sauberer ausgesehen als vorher. Wir
pruefen den Pfad deshalb jetzt einmal beim Submit, statt es nach ein paar hundert
Dateien zu merken. Vielleicht ist ein Befund statt einer Exception fuer dich der
bessere Weg.

\section{Deine acht Punkte}

\subsection*{1. Verlustbehaftete Kompression auf \code{ps}}
\ok\ Du hast recht, und die entscheidende Stelle ist der Zusatz in CF: nicht nur
Koordinaten, sondern auch eine Variable, die von \code{coordinates},
\code{formula\_terms} oder \code{cell\_measures} benannt wird, "even if it is also a
data variable". Genau den hatten wir uebersehen. Im Code stand woertlich:

\begin{quote}
\code{\# ps is a genuine geophysical field and keeps its normal treatment.}
\end{quote}

Das stimmt fuer sich genommen, nur greift die Ausnahme eben trotzdem, weil \code{ps}
zugleich formula\_term der Hybridachse ist. Betraf 9 Dateien. Wir sammeln die
geschuetzten Variablen jetzt nach der Regel statt \code{ps} einzeln auszunehmen.
Nachgemessen war \code{ps} in \code{cli114} tatsaechlich der einzige Treffer.

\subsection*{2. \code{plev}: \code{positive} und \code{units}}
\ok\ Und der Fund war groesser, als er aussieht. Es war kein Fehler an \code{plev}
allein, sondern an allen Vertikalachsen ausser \code{lev}: der Pass, der die
CV-Werte durchsetzt, kannte nur \code{lev}.

\begin{center}
\begin{tabular}{llr}
\toprule
\textbf{Achse} & \textbf{was falsch war} & \textbf{Dateien} \\
\midrule
\code{plev}   & \code{standard\_name} und \code{axis} fehlen, \code{units} "pascal", \code{positive} "up" & 17 \\
\code{height} & \code{axis} fehlt & 35 \\
\code{sdepth} & \code{units} "meter" & 3 \\
\bottomrule
\end{tabular}
\end{center}

Sichtbar wurde es erst, als die Achse richtig \code{plev} hiess. Solange sie
\code{plev19} oder \code{plev3} hiess, lief der Check ins Leere, weil er eine
Variable namens \code{plev} sucht. In \code{cli112} standen also 68 HIGH latent im
Output, ohne dass sie jemand zaehlte.

Dazu der Typ: \code{plev} lag als float32 vor, die Tabelle sagt \code{double}
(17 Dateien), \code{sdepth} als int64. Wird jetzt aus dem \code{type}-Feld der
Tabelle abgeleitet.

\subsection*{3. \code{alevel}: \code{lev}, \code{a}, \code{b} invertieren}
\ok\ \code{stored\_direction} ist "decreasing", die ECMWF-Vertikaltabelle laeuft
top-first, also kam bei uns alles aufsteigend heraus. Gedreht wird die ganze Achse in
einem Zug, nicht nur die Koordinate, damit die Daten an ihrem Level haengen bleiben.
Das ist die einzige Aenderung dieser Runde, die Werte bewegt, deshalb gegen die
echten L137-Koeffizienten geprueft:

\begin{center}
\begin{tabular}{lll}
\toprule
& \textbf{vorher} & \textbf{nachher} \\
\midrule
\code{lev}            & 1{,}00e--05 \dots 0{,}9988 & 0{,}9988 \dots 1{,}00e--05 \\
\code{b}              & 0 \dots 0{,}9988 & 0{,}9988 \dots 0 \\
$p = ap + b \cdot ps$ & steigend & 100880 \dots 1 Pa, streng fallend \\
Testfeld 0\dots136    & 0\dots136 & 136\dots0 \\
\bottomrule
\end{tabular}
\end{center}

Die Bounds-Paare drehen mit, sonst stuende in jeder Zelle weiter [unten, oben].
\code{ap} und \code{b} haben ausserdem das \code{(k)} im \code{long\_name} verloren.

\subsection*{4. \code{deltasigt} mit Absatz im \code{long\_name}}
\ok\ Der Absatz steht schon upstream im \code{long\_name}, rund 900 Zeichen
Methodenbeschreibung mit Literaturangaben. Wir kuerzen jetzt auf den ersten Satz und
verschieben den ganzen Text nach \code{comment}. Das greift generisch, nicht nur bei
\code{deltasigt}.

Gleich mit erledigt: die \code{long\_name} aller Koordinaten kommen jetzt aus der
Tabelle, und zwar ueber den DReq-Namen, nicht ueber den Ausgabenamen. \code{plev19}
heisst "Pressure Levels (19)" und \code{plev3} "Pressure Levels (3)", geschrieben
werden beide als \code{plev}. Bei \code{time} genauso, das sind vier verschiedene
Werte auf derselben Achse. Betraf 492 Dateien.

\subsection*{5. Bodentiefe ohne Bounds, \code{slthick} ohne Gitter}
\ok\ Beides.

Bei \code{sdepth10cm} liefert die Tabelle die Bounds sogar mit
(\code{bounds\_values} "0.0 10.0"), wir haben sie nur nicht gelesen und schrieben
allein den Mittelpunkt 5{,}0 cm. Die Schicht, fuer die der Wert steht, stand also gar
nicht in der Datei.

\code{slthick} kam als \code{slthick(sdepth)} heraus, ein 1D-Feld, dessen Koordinate
der Schichtindex 1..4 war statt einer Tiefe, und im ganzen File kein lat/lon. Liegt
jetzt auf dem Gitter, mit den Schichtmitten als Tiefe und den echten
HTESSEL-Grenzflaechen 0, 7, 28, 100, 289 cm als Bounds. Die aus den Mitten zu
interpolieren waere bei so ungleichen Schichten Unsinn gewesen: die zweite Schicht
ginge dann von 10{,}5 auf 40{,}75 cm statt von 7 auf 28.

\subsection*{6. \code{oplayer4} soll \code{pdepth} heissen}
\ok\ Und dahinter steckte mehr als ein Name. Die drei betroffenen Variablen
\code{scint}, \code{phcint} und \code{absscint} lieferten ein Integral pro
Wassersaeule, verlangt sind vier, eines je Schicht.

Welche vier Schichten das sind, steht in der Tabelle nicht: \code{oplayer4} ist
\code{must\_have\_bounds: yes} mit leerem \code{bounds\_requested}, geliefert werden
nur die vier Werte 15, 50, 136 und 1000 bar. Das laesst sich aber ausrechnen. Euer
Ozeanpapier zum Data Request nennt die Bereiche 0--300, 0--700, 0--2000 m und Rest.
Mit $p = \rho g h$:

\begin{center}
\begin{tabular}{lrrr}
\toprule
& \textbf{15} & \textbf{50} & \textbf{136} \\
\midrule
disjunkt (0--300, 300--700, 700--2000) & 15{,}08 & 50{,}28 & 135{,}75 \\
kumulativ (0--300, 0--700, 0--2000)    & 15{,}08 & 35{,}19 & 100{,}55 \\
\bottomrule
\end{tabular}
\end{center}

Also disjunkte Schichten, drei unabhaengige Treffer auf unter 0{,}3 bar. Der vierte
Wert 1000 gehoert zu einer nach unten offenen Schicht unter 2000 m. Integriert wird
ueber Tiefe, nicht ueber Druck, weil die Schichten in der Quelle Tiefenbereiche sind
und auf der Achse nur als Druck ausgedrueckt werden. Modellschichten, die eine Grenze
kreuzen, werden anteilig geteilt; die Summe der vier ergibt exakt das
Saeulenintegral.

\subsection*{7. \code{lat(lat)} und \code{lon(lon)} auf regulaeren Gittern}
\ok\ Das war der groesste Einzelhebel: 41~Dateien ueber acht Checks, dazu der
Grossteil der \code{AICC006}-Befunde zur Dimensionsreihenfolge.

Die Ursache lag eine Ebene tiefer als der Name. Der DReq benennt eine Dimension nach
der angeforderten Stufe, auf die Platte gehoert der \code{out\_name}, und 110~der
Eintraege in \code{CMIP7\_coordinate.json} unterscheiden sich darin. Wir hatten zwei
Sonderfaelle fest verdrahtet, genau die zwei, die uns aufgefallen waren. Jetzt wird
allgemein aufgeloest.

Darunter lag noch ein zweiter Fehler derselben Sorte: die Liste, die Quelldimensionen
auf DReq-Dimensionen abbildet, war ebenfalls handgepflegt und kannte \code{plev19},
aber nicht \code{plev7h}, und \code{olevel}, aber nicht \code{sdepth}. Eine nicht
zugeordnete Dimension wird nicht umbenannt, sie landet unter ihrem Modellnamen in der
Datei. Deshalb stand bei uns \code{pressure\_levels\_7h} und \code{sdepth} auf der
Platte, und keine \code{out\_name}-Aufloesung haette daran etwas geaendert. Beides
kommt jetzt aus der Tabelle.

\subsection*{8. \code{positive} auf \code{siflcondtop}}
\ok\ Verstanden und umgesetzt: in den Variablentabellen ist \code{positive} eine
Anweisung an den Produzenten, keine Metadatenangabe, und gehoert nicht in die Datei.

Das Vorzeichen selbst haben wir vorher geprueft, bevor wir das Attribut entfernt
haben. \code{siflcondtop} rechnen wir aus \code{ist}, \code{sss} und \code{h\_ice},
und die Formel $k(T_\mathrm{surface} - T_\mathrm{base})/h$ ist bereits die
CMIP-Richtung: im Winter leitet Waerme nach oben aus dem Eis, also negativ. Gemessen
86{,}4\,\% negativ bei einem Mittel von $-24{,}5$\,W\,m$^{-2}$. Die Daten stimmten
also, nur das Attribut musste weg.

\section{Was dein Plugin darueber hinaus gefunden hat}

\begin{center}
\begin{tabular}{lrl}
\toprule
\textbf{Fund} & \textbf{Dateien} & \textbf{Stand} \\
\midrule
\code{coordinates} listet Koordinatenvariablen & 490 & \ok \\
Zeitachse hinter der horizontalen (\code{msftbarot}) & 1 & \ok \\
\code{nz} im \code{coordinates} von \code{tob}/\code{sob} & 2 & \ok \\
\code{tclm} ohne \code{climatology\_bnds} & 1 & \ok, s.\,u. \\
\code{olevel}: \code{lev} fehlt ganz & 6 & \offen, Modellseite \\
zonale Schnitte ohne eigenes Gitter & 4 & \offen \\
\bottomrule
\end{tabular}
\end{center}

\subsection*{\code{pfull} ist jetzt eine echte Klimatologie}

Du hattest recht, dass da eine fehlte. \code{tclm} heisst laut CV zwoelf Werte,
gemittelt ueber eine Anzahl Jahre, mit \code{time: mean within years time: mean over
years}. Wir schrieben die zwoelf Monatsmittel eines einzigen Jahres darunter.

Alles neu zu rechnen scheidet bei uns aus: die Quelle ist 27\,GB pro Jahr bei
194~Jahren, ein Neuaufbau waere rund 5\,TB Lesen und wuerde mit jedem Lauf wachsen.
Wir fuehren deshalb einen Zwischenstand ausserhalb des Laufverzeichnisses, Summe und
Anzahl der Jahre je Monat. Jeder Lauf addiert nur sein eigenes Jahr, die Klimatologie
ist der Quotient und waechst mit. Welche Jahre schon eingerechnet sind, steht mit
dabei, damit ein wiederholt cmorisiertes Jahr nicht doppelt zaehlt.

Das ist bei uns die einzige Datei, die nicht jahrweise archiviert werden kann. Alle
anderen 539 sind es.

\subsection*{Sechs Ozeanvariablen, bei denen die Daten fehlen}

\code{osalttend}, \code{osaltrmadvect}, \code{osaltdiff}, \code{opottemptend},
\code{opottemprmadvect} und \code{opottempdiff} sollen \code{(time, lev, ncells)}
sein und sind bei uns \code{(time, nod2)}.

Das ist kein Metadatenproblem. Die Felder sind im FESOM-Diagnosemodul konstruktiv
zweidimensional, der Quellcode sagt es selbst:

\begin{quote}
\code{real(kind=WP), save, allocatable :: osalttend(:)\ \ ! Salinity tendency - 2D column-integrated}
\end{quote}

Ein XIOS-Eintrag auf ein 3D-Gitter laeuft ins Leere, wenn das Feld nur eine Dimension
hat. Die sechs Regeln sind fuer HR deaktiviert. Der Term pro Level wird im Modell
ohnehin gerechnet und nur aufsummiert, die Aenderung ist also klein; sie geht in die
naechste Modellversion, und die betroffenen Variablen kommen mit dem LR-Lauf.

\section{Offene Fragen an dich}

\subsection*{\code{grid\_label} fuer Dateien ohne horizontales Gitter}

Bei uns sind das 30~Dateien, alle mit \code{-hm-} in der Branding, also globale
Mittel. Sie trugen bisher \code{g130}, \code{g129} oder \code{g122} und behaupteten
damit ein Gitter, auf dem sie nicht liegen.

Wir setzen sie jetzt auf \code{g190}, "A single grid cell for the entire globe". Eine
Zelle fuer die ganze Kugel ist genau das, was ein globales Mittel raeumlich
darstellt, und der Eintrag ist seit Ende Juni gemergt.

Bleibt eine Frage: du hattest geschrieben, die ersten 100 Label seien fuer solche
Faelle reserviert und es sei noch nicht final entschieden. Falls dort etwas anderes
vorgesehen ist, stellen wir um; es sind 29~Regeln in vier Tiers und ein
\code{sed}-Lauf.

\subsection*{Sieben Achsen verlangen Bounds und liefern keine}

\code{oplayer4} ist kein Einzelfall. In \code{CMIP7\_coordinate.json} stehen sieben
Eintraege auf \code{must\_have\_bounds: yes}, haben eine explizite Werteliste und kein
\code{bounds\_requested}:

\begin{quote}
\code{alt16}, \code{alt40}, \code{dbze}, \code{oplayer4}, \code{plev7c},
\code{scatratio}, \code{tau}
\end{quote}

Bei \code{oplayer4} liessen sich die Schichten aus eurem Papier rekonstruieren, wie
oben gezeigt. Fuer die uebrigen duerfte es anderen genauso gehen wie uns.

\subsection*{Ein \code{requested}-Wert, den CF nicht kennt}

\code{vegtype} setzt \code{standard\_name = area\_type}, und der CF-Standardname
\code{area\_type} schreibt vor, dass die Werte aus der CF-Area-Type-Tabelle stammen.
In der \code{requested}-Liste steht aber \code{needleleaf\_evergreen\_tree} im
Singular, den es dort nicht gibt; die anderen drei Baumtypen sind Plural, und die
CF-Tabelle fuehrt durchgehend \code{needleleaf\_evergreen\_trees}. Von 247 geprueften
CF-Begriffen ist das der einzige Ausreisser.

Wir schreiben den CF-Plural, weil derselbe Tabelleneintrag den Standardnamen setzt,
der das verlangt. Dein Plugin meldet das folgerichtig als fehlenden Wert.

\subsection*{Kleinigkeit}

\code{time4} hat im \code{long\_name} einen Tippfehler: "Monthly Mean Daily
Satistics".

\section{Was als naechstes kommt}

\code{cli115} laeuft mit allem aus diesem Dokument. Die Gegenprobe schicken wir nach,
sobald er durch ist.

Offen bleiben von unserer Seite die vier zonalen Schnitte
(\code{msftm}, \code{hfbasin}, \code{sltbasin}), fuer die es kein passendes Gitter im
Register gibt. Die Registrierung ist angestossen. Beim Nachmessen ist uns dabei
aufgefallen, dass unsere drei zonalen Achsen nicht einheitlich sind, einmal 181
Punkte auf ganzen Grad und zweimal auf halben; das ziehen wir vorher gerade, damit
die Registrierung nicht die Uneinheitlichkeit festschreibt.

\end{document}
